% Generated draft: requires mathematical and source review.
\documentclass[11pt]{article}
\usepackage[margin=25mm]{geometry}
\usepackage{fontspec}
\setmainfont{Latin Modern Roman}
\usepackage{amsmath,amssymb,mathtools,unicode-math}
\setmathfont{Latin Modern Math}
\usepackage{enumitem,needspace,xurl,hyperref}
\setlength{\parindent}{0pt}
\setlength{\parskip}{5pt}
\setlength{\emergencystretch}{3em}
\newcommand{\N}{\mathbb{N}}
\newcommand{\Z}{\mathbb{Z}}
\newcommand{\Q}{\mathbb{Q}}
\newcommand{\R}{\mathbb{R}}
\newcommand{\C}{\mathbb{C}}
\newcommand{\F}{\mathbb{F}}
\newcommand{\A}{\mathbb{A}}
\newcommand{\PP}{\mathbb{P}}
\newcommand{\E}{\mathbb{E}}
\newcommand{\eps}{\varepsilon}
\newcommand{\dd}{\,\mathrm{d}}
\begin{document}
\section*{UnsolvedMath ID: 800791}
\textbf{Source problem code:} catalogue entry
\subsection{Definitions and mathematical statement}
Let \(O=0\in\C^n\), and let \(\operatorname{End}(\C^n,O)\) denote the holomorphic map germs \(f:(\C^n,O)\to(\C^n,O)\) fixing \(O\). For this draft, interpret \emph{superattracting} as \(Df(O)=0\), and define the order of a nonzero germ by its first nonzero homogeneous term:
\[
f(z)=\sum_{k\ge d}F_k(z),\qquad F_d\not\equiv0,\qquad \operatorname{ord}(f)=d,
\]
where each \(F_k:\C^n\to\C^n\) is a homogeneous polynomial map of total degree \(k\).

Two such germs \(f\) and \(g\) in the same dimension are \emph{formally conjugate} if there exists
\[
H\in (z_1,\ldots,z_n)\C[[z_1,\ldots,z_n]]^n,
\qquad \det DH(O)\ne0,
\]
satisfying \(H\circ f=g\circ H\) as formal power series. Thus the formal coordinate change is invertible; its linear part is allowed to be any invertible complex matrix.

\textbf{Problem OP10 (Abate, 2011).} Determine the formal conjugacy classes of superattracting germs in each of the following cases:
\[
\begin{aligned}
&f\in\operatorname{End}(\C^n,O),\quad \operatorname{ord}(f)=2,
&&\text{for every integer }n\ge3;\\
&f\in\operatorname{End}(\C^2,O),\quad \operatorname{ord}(f)=3.
\end{aligned}
\]
The supplied excerpt imposes no additional nondegeneracy or finiteness hypothesis on the leading homogeneous term. It does not supply the surrounding definitions of superattracting, order, or formal conjugation; the conventions made explicit above require confirmation against the cited source.
\subsection{Short English statement}
Abate's Problem OP10 asks for a classification under invertible formal changes of coordinates of superattracting holomorphic germs of order \(2\) in every complex dimension \(n\ge3\), and of order \(3\) in complex dimension \(2\). The draft states explicit conventions for the terminology, which require source review. The registry labels the problem open; that label has not been independently verified.
\subsection{Sources}
\begin{itemize}
\item Supplied local UnsolvedMath registry record: permanent ID \(800791\), source problem code \texttt{catalogue entry}, title ``Local Discrete Holomorphic Dynamics --- Open Problem OP10.'' The record attributes the problem to Abate, with proposed year \(2011\), and supplies the status \texttt{open}; this is registry metadata, not an independently verified literature assessment.
\item Abate, \emph{Open Problems in Local Discrete Holomorphic Dynamics} (2011), Problem OP10. Supplied source URL: \url{https://arxiv.org/abs/1107.4876}. The supplied provenance records access on \texttt{2026-10-16} and extraction by \texttt{arxiv\_latex}. These are supplied provenance statements; no independent source visit or verification is claimed.
\end{itemize}
\end{document}
